\section{Measuring the two noise scales from a single batch}
\label{sec:estimation}

Both $\tau_b$ and $\tau_w$ are recoverable from gradients a trainer already computes. Index the
microbatches by prompt block and by rollout sub-group and accumulate them separately:
\begin{equation}
  \label{eq:cells}
  g_{ab} := \text{mean gradient over block } a \text{ using sub-group } b,
\end{equation}
for $a = 1 \dots K$ and $b \in \{1,2\}$, each computed with its own sub-group's advantages over
$P/K$ prompts and $G/2$ responses.
Blocks hold disjoint prompts, and given a prompt the two sub-groups are independent, so every term
below is an inner product of independent quantities:
\begin{align}
  \label{eq:estimators}
  \gbar^{\!\top} M \gbar &= \E\big[ g_{a1}^{\!\top} M\, g_{a'2} \big]_{a \neq a'}, \\
  \tr(M\Sb) &= \tfrac{P}{K}\Big( \E\big[ g_{a1}^{\!\top} M\, g_{a2} \big]
     - \gbar^{\!\top} M \gbar \Big), \notag \\
  \tr(M\Sw) &= \tfrac{P}{2K}\, \E\big[ \| g_{a1} - g_{a2} \|_M^2 \big]. \notag
\end{align}
The first pairs different prompts with different responses, the second the same prompts with
different responses, and the third cancels the prompt-level term inside the difference.

The cost is $2K$ gradient buffers rather than one, no extra responses, and no extra backward work
beyond tagging each microbatch -- which any trainer using gradient accumulation is already
positioned to do. Raising $K$ averages the block statistics over $K(K-1)$ ordered pairs rather than
one and shrinks the $P/K$ multiplier in front of the noisiest term, so the estimate improves with
the number of microbatches the trainer already uses. $M = I$ gives the practical form; $M = F$ is
obtained from measured KL rather than by materialising the Fisher.

Getting this right requires handling three biases, each easy to walk into.

\paragraph{The subtractive form is badly conditioned.} Estimating the total variance and
subtracting the within-prompt part differences two large, highly correlated quantities. On
homogeneous pools it returns a negative $\tau_b$ and an infinite $G^{\star}$; where it returns something
finite, it compresses the answer toward the middle of the range. Measured against exactly known
values over pools whose noise ratio spans a factor of $26$, \eqref{eq:estimators} recovers $G^{\star}$
to a median error of $1.3\%$ and a worst case of $6.2\%$, while the subtractive form reports
$6.1$ where the truth is $9.4$ and $6.9$ where the truth is $15.2$.

\paragraph{A repeated prompt lands in two blocks.} Sampling prompts with replacement puts the same
prompt in two blocks, so the different-block inner product picks up a same-prompt term and
$\tau_b$ is biased down -- by $45\%$ on the most homogeneous pool tested. Drawing distinct prompts
within a batch, which iterating a dataset does anyway, removes it.

\paragraph{A finite corpus makes blocks anticorrelated.} Drawing without replacement from a corpus
of $N$ has the opposite effect on the other term: two blocks hold distinct prompts, so
\begin{equation}
  \label{eq:finite-corpus}
  \E\big[ g_{a1}^{\!\top} M\, g_{a'2} \big]_{a \neq a'}
  \;=\; \gbar^{\!\top} M \gbar \;-\; \frac{\tr(M\Sb)}{N-1},
\end{equation}
which is exact and which the estimator corrects using the $\tau_b$ it already has. The correction
matters whenever $\tau_b$ is large relative to the signal -- that is, whenever the run is far below
its critical batch size, which is the usual case. Uncorrected, the signal estimate on one pool came
out negative. The same factor appears in the theory: with a finite corpus the between-prompt
variance enters $\Cov(\ghat)$ multiplied by $1 - (P-1)/(N-1)$, so
\begin{equation}
  \label{eq:gstar-corpus}
  G^{\star} \;=\; 1 + \sqrt{\frac{\tau_w}{\big(1 - \tfrac{P-1}{N-1}\big)\tau_b}},
\end{equation}
and a run that covers its whole corpus every step should use larger groups than one that samples
from a large corpus.

\paragraph{Preconditioners.} An update is $\eta\,\Pi\,\ghat$ for a preconditioner $\Pi$, and Adam
is scale invariant, so both the drift model and the estimators must be read in the metric the
update actually moves in. In practice this means instrumenting the preconditioned gradient: the
results carry over with $F \mapsto \Pi F \Pi$ under the standard assumption that $\Pi$ varies
slowly relative to batch-to-batch fluctuation.

\paragraph{Controlling drift.} Since $D \propto \eta^2$, a multiplicative controller
\begin{equation}
  \eta_{t+1} \;=\; \eta_t \cdot \mathrm{clip}\Big( \big( D^{*} / D_t \big)^{1/2},\ c^{-1},\ c \Big)
\end{equation}
corrects a mis-set step in a single update when the drift estimate is clean. Measured over runs on
transformers, it holds the median drift at its target with no detectable bias in log space.
Exponential smoothing of $D_t$, the obvious refinement, makes it worse: the lag biases the realised
drift low by roughly a factor of two.
